\chapter{Croup}

\section*{Definition}
Croup (acute laryngotracheobronchitis) is a common respiratory illness in children characterized by a bark-like cough, inspiratory stridor, and hoarseness due to subglottic airway edema. The Westley croup score (0--17) stratifies severity based on stridor, retractions, air entry, cyanosis, and level of consciousness.

\section*{What Happens in Your Body}
Viral infection of the laryngeal and tracheal mucosa causes edema and inflammation in the subglottic region (the narrowest part of the pediatric airway). This produces dynamic airway narrowing during inspiration, creating the characteristic inspiratory stridor and seal-like barking cough.

\section*{Causes}
\begin{itemize}
  \item \textbf{Viral (most common):} Parainfluenza virus types 1, 2, and 3 (75\% of cases), RSV, influenza A and B, adenovirus, metapneumovirus
  \item \textbf{Spasmodic croup:} Recurrent, non-infectious, possibly allergic or reflux-mediated, presents suddenly at night
  \item \textbf{Bacterial (rare):} \textit{Corynebacterium diphtheriae} (diphtheritic croup), \textit{Mycoplasma pneumoniae}
  \item \textbf{Other:} GERD-induced laryngeal edema mimics croup symptoms in some cases
\end{itemize}

\section*{When to See a Doctor}
See your doctor if:
\begin{itemize}
  \item Stridor at rest (not just with agitation or crying)
  \item Severe retractions (supraclavicular, intercostal, substernal), nasal flaring, or head bobbing
  \item Oxygen saturation less than 92\% or cyanosis
  \item Altered mental status, irritability, or lethargy
  \item Inability to swallow or drooling (differentiate from epiglottitis)
\end{itemize}

\section*{Related Symptoms}
\begin{itemize}
  \item Barking cough, inspiratory stridor, hoarseness, low-grade fever, coryza, worse at night
\end{itemize}
\textbf{Important:} Croup typically affects children aged 6 months to 3 years and symptoms peak on nights 2--3; differentiate from epiglottitis (acute onset, high fever, drooling, tripod positioning) which requires urgent airway management.

\section*{Self-Care / Home Management}
\begin{itemize}
  \item Calm the child (crying worsens stridor and airway obstruction)
  \item Cool mist humidifier or steam from a hot shower (humidity may soothe the airway)
  \item Acetaminophen or ibuprofen for fever and discomfort
  \item Hydration with clear liquids in small, frequent amounts
  \item Elevate the head during sleep to reduce supine airway edema
\end{itemize}

\section*{How Your Doctor Will Evaluate You}
\begin{itemize}
  \item Clinical diagnosis based on characteristic barky cough and stridor
  \item Westley croup severity score for objective assessment: mild (score less than 3), moderate (3--6), severe (greater than 6)
  \item Neck X-ray (AP view) may show subglottic narrowing (steeple sign) if diagnosis is uncertain
  \item Pulse oximetry to assess oxygenation in moderate-to-severe cases
  \item Differentiate from epiglottitis (lateral neck X-ray shows thumb sign), bacterial tracheitis, or foreign body aspiration
\end{itemize}

\section*{Treatment Options}
\begin{itemize}
  \item Single dose of oral or intramuscular dexamethasone (0.6 mg/kg, max 10 mg) for all severity levels
  \item Nebulized epinephrine (racemic epinephrine 0.25--0.5 mL or L-epinephrine 1:1000, 0.5 mL/kg) for moderate-to-severe croup -- observe for 2--4 hours for rebound
  \item Heliox (helium-oxygen mixture) for severe croup to reduce turbulent airflow
  \item Hospital admission for persistent stridor at rest, hypoxia, or recurrent need for nebulized epinephrine
  \item Intubation for impending respiratory failure (rare, less than 1\% of cases)
\end{itemize}

\clearpage
